\section{Análisis de variables y utilidad para modelamiento}

La tabla siguiente integra descripción estadística, relación observada y decisión técnica. Una
correlación débil no implica inutilidad multivariable: interacciones y relaciones no lineales pueden
aportar información, pero deben comprobarse fuera de muestra.

\small
\begin{longtable}{p{3.2cm}p{3.4cm}p{3.2cm}p{4.2cm}}
\caption{Síntesis de variables para la siguiente etapa.}\label{tab:analisis-variables}\\
\toprule
Variable & Descripción estadística & Relación con popularidad & Utilidad o riesgo \\
\midrule
\endfirsthead
\toprule
Variable & Descripción estadística & Relación con popularidad & Utilidad o riesgo \\
\midrule
\endhead
\code{danceability} & Continua 0--1, distribución amplia & Pearson 0,064; Spearman 0,055; curva por cuantiles & Candidata; permitir no linealidad \\
\code{energy} & Continua 0--1 & Asociación casi nula con target & Candidata, pero redundante con sonoridad y acousticness \\
\code{loudness} & dB con cola negativa & Pearson 0,072; Spearman 0,067 & Candidata; escalar y revisar extremos \\
\code{speechiness} & Masa cerca de cero y cola alta & Pearson -0,047; Spearman -0,067 & Candidata; posible transformación y segmentación \\
\code{acousticness} & Continua 0--1, relación no lineal & Pearson -0,039; Spearman 0,010 & Candidata con cautela por redundancia \\
\code{instrumentalness} & Acumulación en cero y valores altos & Mayor asociación: -0,128/-0,124 & Candidata; no convertir asociación en regla \\
\code{liveness} & Asimétrica y acotada & Cerca de cero & Mantener para prueba multivariable \\
\code{valence} & Distribución amplia 0--1 & Cerca de cero & Candidata; relacionada con danceability \\
\code{tempo} & Continua; incluye cero y máximo 243,372 & Cerca de cero & Validar extremos y escalar \\
\code{duration_ms} & Cola larga; máximo 87,29 minutos & Pearson -0,023; Spearman 0,013 & Evaluar \code{log1p} y escala robusta \\
\code{explicit} & 7.704 verdaderas de 89.740 & Diferencia media 4,03; $d=0{,}196$ & Binaria candidata; auditar por confusores \\
\code{track_genres} & 114 etiquetas básicas; combinaciones multietiqueta & Diferencias descriptivas por género & Multi-hot; evitar doble conteo \\
\bottomrule
\end{longtable}
\normalsize

\subsection{Cardinalidad de variables categóricas}

\begin{table}[H]
\centering
\caption{Cardinalidad sobre el dataset procesado.}
\label{tab:cardinalidad}
\resizebox{\textwidth}{!}{\begin{tabular}{lrrl}
\toprule
 & Valores unicos & Porcentaje sobre filas & Riesgo \\
variable &  &  &  \\
\midrule
artists & 31437 & 35.031 & Alta cardinalidad y memorizacion \\
album\_name & 46589 & 51.916 & Alta cardinalidad y memorizacion \\
track\_name & 73608 & 82.024 & Alta cardinalidad y texto libre \\
track\_genres & 1432 & 1.596 & Multietiqueta; combinaciones de etiquetas \\
explicit & 2 & 0.002 & Desbalance y posibles variables de confusion \\
key & 12 & 0.013 & Categoria nominal codificada como numero \\
mode & 2 & 0.002 & Baja cardinalidad \\
time\_signature & 5 & 0.006 & Categoria discreta con clases poco frecuentes \\
\bottomrule
\end{tabular}
}
\end{table}

\code{artists}, \code{album_name} y \code{track_name} podrían memorizar identidades o ejemplos
concretos. El dato de 1.432 valores únicos de \code{track_genres} cuenta combinaciones de etiquetas,
no géneros básicos. \code{key}, \code{mode} y \code{time_signature} están almacenadas como números,
pero su significado es categórico; tratarlas como distancias continuas impondría relaciones
arbitrarias.

\subsection{Multicolinealidad y redundancia}

La multicolinealidad ocurre cuando predictores aportan información lineal muy similar. No afecta
necesariamente la predicción global, pero puede aumentar la varianza de coeficientes, cambiar signos
y dificultar explicaciones. Aquí, \code{energy}, \code{loudness} y \code{acousticness} forman el
principal grupo de atención.

El escalamiento coloca variables en rangos comparables y es necesario para penalizaciones, aunque
no elimina correlación. Ridge puede estabilizar coeficientes distribuyendo señal entre predictores;
Lasso puede seleccionar, pero su elección entre variables correlacionadas puede ser inestable. El
factor de inflación de varianza (VIF) sería útil si se adopta una regresión lineal explicativa; no se
calcula ahora porque aún no se definió la matriz final de diseño ni el propósito inferencial.
